% !TEX root = ../metatime_monograph.tex
\chapter{The Metatime Framework}
\label{ch:metatime_framework}

% CITATION_CONTEXT_V10_9_BEGIN
\paragraph{Established context.}
The standard geometric ingredients behind phase holonomy, quantum-state metrics,
and cyclic evolution are described by Pancharatnam, Simon, Berry,
Aharonov--Anandan, and Provost--Vall\'ee
\cite{Pancharatnam1956,Simon1983,berry1984,AharonovAnandan1987,ProvostVallee1980}.
Connections to information geometry are conventional in the sense of Fisher,
Rao, and Amari--Nagaoka \cite{Fisher1925,Rao1945,AmariNagaoka2000}; the specific
Metatime synthesis and its operator assignments are the construction tested in
this monograph.
% CITATION_CONTEXT_V10_9_END

\section{The Information Constant \(\kappa\)}

The Metatime framework defines the information constant
\begin{equation}
\kappa\equiv\frac{\ln2}{24\pi}.
\label{eq:kappa_construction}
\end{equation}
Its numerical value is
\[
\kappa=0.009193150006360484\ldots .
\]

\claimstatus{B -- model postulate with structural motivation; the arithmetic identities below are exact, while the physical assignments of their factors are model-level.}

This claim status is the one used consistently in Appendix~\ref{app:kappa}.  The
dedicated information-spinor analysis under
\path{TIR/zeta_information_axis/}
and the crosswalk in Appendix~\ref{app:information_spinor_crosswalk} sharpen the
internal factorization of the normalization without promoting it to an
established first-principles invariant of quantum field theory or differential
geometry.

\subsection{Binary information}

For a balanced binary distribution, Shannon entropy in natural units is
\cite{shannon1948}
\begin{equation}
H_{\rm bin}\!\left(\frac12\right)=\ln2.
\label{eq:shannon_binary}
\end{equation}
This is an established information-theoretic identity.  It supplies the
numerator of \eqref{eq:kappa_construction}; it does not by itself determine the
denominator.

\subsection{Normalized recurrence before radians}

A cyclic process may first be represented by a normalized turn coordinate
\begin{equation}
q\in\mathbb R/\mathbb Z,
\end{equation}
so that one recurrence is $q\mapsto q+1$.  The intrinsic frequency is a winding
rate, $f_q=\Delta N/\Delta t$, before an angular convention is selected.  An
angular coordinate is introduced only by
\begin{equation}
\phi=Cq,
\end{equation}
where $C$ is the closure period in the chosen angular unit; in radians
$C=2\pi$.

For the two-level state used in the information-spinor subprogramme, the
balanced equatorial loop carries Berry holonomy $-1$.  This is a standard
geometric-phase statement for the specified state family
\cite{Simon1983,berry1984}.  It should not be replaced by the looser claim that
an arbitrary ``complete cycle'' has solid angle $4\pi$: Berry phase depends on
the actual closed path and its enclosed solid angle.

\subsection{Discrete cross-factorization}

The exact arithmetic identity
\begin{equation}
24=8\cdot3=12\cdot2=6\cdot4
\label{eq:twentyfour_crossfactor}
\end{equation}
provides several internally compatible decompositions.  Their Metatime/TIR
interpretations are model assignments:
\begin{itemize}
\item $8\cdot3$: eight declared mixing sectors times three flavours;
\item $12\cdot2$: twelve projective cycles times two half-turn units;
\item $6\cdot4$: six spinor cycles times four half-turn units.
\end{itemize}
The equality of the integers is exact; the physical interpretation of the
factors is not an independent theorem.  The older tetrahedral notation remains
a useful structural representation: $|A_4|=12$ for rotational tetrahedral
symmetry and $|S_4|=24$ for the full tetrahedral symmetry.  The use of these
groups is part of the NOEMA/Metatime model geometry and must not be conflated
with the Weyl group of $SU(3)$, which is $S_3$.

\subsection{Conditional arithmetic--geometry reconstruction of \(\kappa\)}

If the Metatime information cycle assigns one binary information quantum to
twelve normalized projective cycles,
\begin{equation}
\frac{\dd I}{\dd q}=\frac{\ln2}{12},
\label{eq:info-per-projective-turn}
\end{equation}
then conversion to an angular coordinate with closure period $C$ gives exactly
\begin{equation}
\frac{\dd I}{\dd\phi}=\frac{\ln2}{12C}.
\end{equation}
For the radian representation $C=2\pi$,
\begin{equation}
\frac{\dd I}{\dd\phi}=\frac{\ln2}{24\pi}=\kappa.
\end{equation}
Thus the present formulation separates the model's arithmetic cycle assignment
from geometric closure.  It is a useful cross-derivation and implementation
constraint, but it does not remove the model status of the twelve-cycle
assignment itself.  Appendix~\ref{app:kappa} remains the authoritative
publication claim boundary.

\section{The Exponential Mass Formula}

The central postulate of the Metatime framework is that every physical mass
is determined by an action \(S\) on the Poincar\'{e} disk via an exponential
Boltzmann-like factor:

\begin{equation}
m = E_P \, e^{-S/\kappa}.
\label{eq:mass_formula}
\end{equation}

Here \(E_P = \sqrt{\hbar c^5 / G} = 1.22089 \times 10^{22}\)~MeV is the Planck
energy (the fundamental scale of quantum gravity), and \(S\) is a
dimensionless action constructed from the L-constants and quark primes.

The exponential form is a Metatime modelling ansatz motivated by geometric
phase accumulation and information-weighted state counting.  It is not derived
from the standard canonical partition function by a theorem identifying
$\kappa$ with a thermodynamic temperature.

\subsection{Why the Exponential Form?}

For comparison, a canonical statistical system uses
\[
Z = \sum_i e^{-E_i / k_B T}.
\]
The Metatime construction adopts an analogous exponential architecture with
$S/\kappa$ dimensionless.  The analogy motivates the functional form; it does
not identify $S$ with ordinary thermodynamic entropy or $\kappa$ with $k_BT$.
Physical relevance is therefore assessed by the explicit numerical and
prospective tests in later chapters.

\section{Action Construction Principles}

Each particle's action \(S\) is built from the following elements:

\begin{enumerate}
\item \textbf{Spin degeneracy:} A term \(\frac{1}{2}\) for spin-1/2 fermions,
encoding the two spin states.
\item \textbf{Generation terms:} Terms proportional to \(\kappa\) for each
generation, representing the information deficit shared across generations.
\item \textbf{Flavour terms:} Terms involving quark primes \(q_i\) and
L-constants \(L_3, L_4, L_5\), encoding the flavour structure.
\item \textbf{Geometric corrections:} Terms involving ratios of L-constants
and powers of \(\kappa\), representing higher-order curvature corrections on
the Poincar\'{e} disk.
\end{enumerate}

\section{The NOEMA Tetrahedron}

The NOEMA (Neutral Operator Eigenvalue Mapping Algebra) tetrahedron is a
geometric construction in the three-dimensional operator space spanned by:

\[
\{\hat I,\ \hat M,\ \hat A\},
\]

where \(\hat I\) is the isospin operator, \(\hat M\) is the mass operator,
and \(\hat A\) is the axial operator.  The tetrahedron encodes the mixing
patterns of quarks and leptons:

\begin{itemize}
\item The edge lengths are determined by \(L_3\) and \(L_4\).
\item The face areas determine mass splittings within the model.
\item The dihedral angles supply candidate mixing-angle coordinates.
\end{itemize}

The NOEMA tetrahedron provides the model geometry used for the CKM and PMNS
constructions discussed in Chapters~\ref{ch:ckm_matrix} and
\ref{ch:pmns_mixing}.  Those phenomenological assignments are distinct from
the exact group orders appearing in \eqref{eq:twentyfour_crossfactor}.

\section{External Scales and Inputs}

The framework requires two external energy scales and three external numerical
inputs:

\begin{description}
\item[\(E_P = 1.22089 \times 10^{22}\)~MeV:] The Planck energy, setting the
absolute mass scale via the exponential formula.
\item[\(E_{\text{proton}} = 938.272\)~MeV:] The proton mass, used as the
reference scale for baryonic and electroweak quantities.
\item[\(N_c = 3\):] The number of colours in QCD, needed for hypercharge
assignments and anomaly cancellation.
\item[\(Y(L) = -1/2\):] The lepton doublet hypercharge, required for anomaly
cancellation checks.
\item[\(\text{Crewther factor} = 2.4 \times 10^{-16}\):] The chiral
perturbation theory factor relating \(\theta_{\text{QCD}}\) to the neutron
electric dipole moment.
\end{description}

The remaining numerical relations are generated from the declared structural
inputs of the model.  The phrase ``no continuous fit'' should not be read as
``no assumptions'': discrete assignments, external scales, and retrospective
operator choices are tracked separately by the publication claim hierarchy.
